\chapter{Acceptance and implementation sequence}
\label{ch:validation}

\section{Three different questions require different evidence}

Engineering correctness asks whether the implementation follows the specified
types, transitions and arithmetic. Accounting validity asks whether the
selected rules faithfully implement the applicable standards and disclosed
policies. Scientific validity asks whether the economic assumptions and
uncertainty representation support the intended forecasts or comparisons.
A balanced journal answers only part of the first question.

The implementation must therefore maintain separate validation results.
Compilation and deterministic examples cannot establish forecasting quality.
Agreement with a public report cannot identify an undisclosed contract term.
A lower extraction error on a document benchmark cannot establish that the
engine preserved every materially different scenario response.

\section{Core conformance cases}

\begin{longtable}{L{16mm}L{73mm}L{47mm}}
\toprule Test & Required case & Pass condition\\\midrule\endhead
T01 & Exact debit/credit batches, including reversals and period closing &
Every committed book conserves signed balances; imbalance rejected atomically.\\
T02 & Money, quantities, currencies, instants and interval flows &
Invalid additions and context combinations rejected; explicit conversions traced.\\
T03 & Original filing, later correction and as-known cutoff &
Earlier query reproduces the original evidence view; later view shows revision.\\
T04 & Annual and interim cumulative flows with revised comparatives &
Quarter derivation only across compatible contexts and vintages.\\
T05 & Fixed/floating borrowing with identical historical postings &
Both explanations survive; next-year expense is 5 and 7 under the stated shock.\\
T06 & Continuous fixed share $\alpha\in[0,1]$ &
Exact model returns $I=7-2\alpha$ and extrema 5, 7 without assigning probabilities.\\
T07 & Duplicated encoding of one history &
Canonical duplicates preserve results and declared weights; no extra
probability mass from an additional encoding.\\
T08 & Linked interpretation across debt classification and payment treatment &
Impossible combinations rejected through shared constraints.\\
T09 & Rounded totals, duplicate facts and contradictory notes &
Correct interval compatibility; no automatic zero or balancing plug.\\
T10 & Infeasible or unsupported constraint problem &
Conflicting sources exposed; timeout/unsupported case reported as unknown.\\
T11 & Zero-coupon bond and ordinary equity examples &
Carrying amounts, payments, share quantities and report mappings match derivations.\\
T12 & Foreign-currency asset, translation and cash movement &
No mixed-unit balance test; translation effects separated from cash flows.\\
T13 & Acquisition and intragroup balance elimination &
Perimeter and adjustment book reproduce group totals without inventing capex.\\
T14 & Conflicting recognition rules and unknown guards &
Explicit conflict or symbolic partition; file order has no effect.\\
T15 & Insufficient cash with constrained funding &
Funding failure/default retained as an outcome; no hidden borrowing.\\
T16 & Expected-loss measurement before future default is realised &
Measurement uses only information allowed at that date.\\
T17 & Repeated/reordered independent events and interrupted execution &
Idempotent replay; branch isolation; identical committed outputs.\\
T18 & Conditional distributions and optional weighted mixture &
Branch distributions retained; mixture agrees with declared weights and total
probability; unweighted family never labelled probabilistic.\\
\bottomrule
\end{longtable}

For synthetic cases, expected histories, excluded histories and scenario
consequences are known by construction. Tests should include counterexamples
that balance but have the wrong economic effect, such as treating a
non-cash acquisition movement as operating cash. Such tests discriminate
between arithmetic validity and the intended semantics.

Differential tests against an existing ledger check only the shared supported
subset. ACTUS schedule comparisons check contractual events, and QuantLib
comparisons check selected valuation calculations under matched conventions.
Disagreements must be classified as different conventions, unsupported cases
or implementation errors before one tool is treated as an oracle.

\section{Extraction and reconstruction evaluation}

A company corpus must include different standards, industries, languages,
filing formats, extensions, interim calendars and restatements. Company and
time splits are required so repeated report templates or revised versions do
not leak into evaluation. Public benchmark tables are useful development
material, while acceptance uses complete disclosure packages that include
notes and their relationships.

Reviewers annotate source facts, contexts, mappings, required relationships,
known ambiguities and unsupported assertions. Relevant measures include
correct source-linked values and contexts, unresolved extraction volume,
coverage of note obligations, reconciliation failures and invented facts.
Report each separately. A single overall score can hide a serious error in a
small but consequential financing note.

Historical reconstruction should test whether the known synthetic history
remains feasible, whether known incompatible histories are rejected, and
whether economically distinct alternatives are preserved. On real companies,
public data generally cannot establish the true journal. Review instead
assesses evidential compatibility, unsupported commitments and conditional
economic consequences. Later disclosures may provide discriminating evidence
for previously ambiguous terms.

No universal numerical threshold for semantic coverage or extraction accuracy
is fixed in this document. Before a release evaluation, an accounting-reviewed
test corpus and task-specific acceptance thresholds must be recorded in an
evidence contract. Hard conformance failures, fabricated citations, silent
discarding of valid alternatives and undocumented accounting adjustments veto
acceptance regardless of an aggregate score.

\section{Evaluate the DSL against the strongest practical baseline}

The comparison uses the same IR, accounting examples and tasks. One group of
models is authored through the typed Python API; another through the DSL.
Both receive equivalent documentation and examples. Tasks include encoding a
movement schedule, adding an alternative interpretation, changing an effective
date and tracing a simulated amount. Measure correctness, material omissions,
time to a reviewed result, ease of explaining a change and maintenance effort.

This comparison tests whether the notation and compiler help users; it is not
a comparison between careful DSL code and unstructured spreadsheets. If the
DSL offers little advantage, retain the shared IR and library and narrow the
textual language. Small pilots nominate improvements. A claim of superiority
requires a predeclared comparison and uncertainty evidence, not a few favourable
demonstrations.

\section{Forecasting evaluation comes later}

Forecast evaluation uses historical information cutoffs, fixed data vintages,
declared economic scenarios and proper separation of training and testing
periods. Baselines should include a naive persistence or ratio model, a tuned
conventional financial model, the structural model with a single selected
interpretation, and the same model preserving alternatives. This identifies
whether any benefit comes from better economics, more information or the
representation of ambiguity.

Construct the inexpensive comparators before choosing a complex forecasting
method. A no-growth operating forecast with scheduled contractual financing
tests whether the model adds value beyond existing obligations. A constant
margin and turnover model tests whether detailed operations improve on familiar
financial ratios. A trailing-average growth model tests whether fitted dynamics
improve on recent experience. A fixed-exposure contractual stress model tests
whether estimated management responses help beyond repricing the existing
book. Each comparator must generate coherent accounts through the same reporting
rules and use the same historical information. In synthetic cases, the known
hidden state also supplies an oracle conditional on the stipulated scenario.

Evaluate these comparisons separately under stable conditions, large rate
resets, refinancing pressure, credit deterioration and changes in reporting
perimeter. A pooled average can conceal failure in the situation the product
is intended to explain. Supported underperformance against a cheap comparator
in a declared critical situation vetoes promotion of the complex forecasting
method and must lead the result discussion. These checks are falsification
tests, not training targets or evidence of superiority when merely passed.

Evaluate point forecasts where they are a stated target, distributional
calibration when probabilities are justified, and decision-relevant robustness
when only feasible sets are available. A wide feasible set is not automatically
a good prediction, and an unweighted set cannot be assessed as if it were a
95\% predictive interval. In multi-scenario stress tests, scenario severity and
probability remain separate inputs.

Any serious numerical comparison must follow the repository's research plan,
assumption audit, baseline, veto, uncertainty and result-record policies. The
present document runs only deterministic checks of its own examples. It
neither selects a scientific forecasting default nor reports comparative
performance.

\section{Implementation milestones and dependency order}

\begin{longtable}{L{16mm}L{76mm}L{44mm}}
\toprule Stage & Concrete implementation deliverable & Exit evidence\\\midrule\endhead
P0 & Evidence registry, contextual types, versioned IR and exact posting engine;
Python constructors for the running cases & T01--T04, T09 and T17; source
inspection and replay.\\
P1 & Finite choices, shared symbols, linear solver adapter, historical bridges,
conditional simulation and report explanations & T05--T10, T14 and T18;
synthetic family preserved.\\
P2 & Textual Core 0 frontend with identical IR; initial corporate package;
document/XBRL import adapters & Parser/API equivalence; T11--T13; reviewed
three-date public-company case.\\
P3 & Contract and valuation adapters, dated standard packages, richer financing
and tax/funding policies & Contract/valuation differential cases; accounting
review; T15 and transition cases.\\
P4 & Bank package, allowance models, regulatory reconciliation and management
response alternatives & T16; banking-specific reviewed histories and stresses.\\
P5 & Insurance cohort package, reinsurance and actuarial measurement adapters &
Reviewed cohort roll-forwards, information-date and service/finance tests.\\
P6 & Optional probabilistic inference, model-weight sensitivity and broader
scenario research & Explicit targets, sampler/numerical validity, prior
sensitivity and downstream evaluation.\\\bottomrule
\end{longtable}

These stages specify dependency order, not dates or staffing estimates.
Banking and insurance representation requirements influence P0 types, although
their measurement engines arrive later. Probability weighting is optional;
P1 already produces useful conditional cases and bounds. A failure of a chosen
extractor or inference algorithm triggers a local replacement or repair,
rather than abandoning the product's representation requirements.

\section{Decisions and replaceable assumptions}

The requirements fixed by this design are preservation of evidence and
alternatives, explicit time and basis, balanced and admissible posting, economic
state beyond journals, and distinct set/probability semantics. Replaceable
choices include the storage engines, parser generator, solver, numerical
libraries, aggregation resolution and user interface. Every scientifically
material aggregation, behavioural parameter, prior or truncation rule must
record provenance, purpose, failure mode and a diagnostic before it becomes a
default.

The immediate next implementation task is P0 with the synthetic financing
cases and their evidence records. This provides a small testable foundation
for the entire design. Broader corporate coverage should grow through reviewed
examples, not through a claim that one generic rule package represents all
financial statements.
